\chapter{Minería de datos: encontrar patrones}
\label{cap:mineria}
\textit{Criterio 3 de la rúbrica (12\,\%): identifica patrones, tendencias y relaciones con técnicas pertinentes y los
interpreta en función del problema.}

\section{El índice de presión turística (IPT)}
\label{sec:ipt}
\textbf{Pregunta.} ¿Qué lugar está más presionado, si cada uno se mide con cosas distintas? Chetumal tiene Tren Maya y
cruces de Belice; la Ruta tiene visitantes a zonas arqueológicas; Cancún tiene ocupación hotelera semanal. Hace falta una
sola escala de 0 a 1.

\begin{intuicion}
Es como calificar a alumnos de grupos distintos con exámenes distintos. Primero se pasa cada examen a la misma escala
(de 0 a 1, donde 0 es la peor calificación de todo el estado y 1 la mejor) y luego se promedian los exámenes que cada
alumno presentó. Al final se dice quién está en la mitad baja, quién en el 10\,\% más alto y quién en medio.
\end{intuicion}

\subsection{Los datos: el panel mensual}
Una tabla de 15 lugares $\times$ 55 meses (enero de 2022 a julio de 2026) = 825 renglones, con: personas que bajan del
Tren Maya, cruceristas y visitantes INAH \textbf{por cada mil habitantes}, llegadas por cuarto de hotel y ocupación. Los
diez destinos del resto del estado dan la escala común: sin ellos, ``tranquilo'' no significaría nada.

\textbf{Un hallazgo al armarlo}: la serie ``zonas arqueológicas'' de SITUR-Q es el mismo dato del INAH sumado por destino
(Chetumal 2025: 39,459 = Oxtankah 11,017 + Kohunlich 21,850 + Dzibanché 6,592). Usar las dos habría contado dos veces a
los mismos visitantes; se usa solo el INAH, zona por zona.

\subsection{Las fórmulas y su fundamento}
\textbf{Paso 1, llevar todo a la misma escala (mín–máx).}
\[
z_{k,d,t}=\frac{x_{k,d,t}-\min_k}{\max_k-\min_k},
\]
donde $\min_k$ y $\max_k$ son el mínimo y el máximo del componente $k$ \textbf{en todos los lugares y meses juntos}.
\emph{Fundamento}: la transformación es lineal, así que conserva el orden y las proporciones; el 0 y el 1 tienen
significado (``lo más bajo'' y ``lo más alto registrado en el estado''). \emph{Alternativa}: estandarizar
($z=(x-\bar x)/\sigma$) da números sin límite y negativos, más difíciles de explicar; usar percentiles pierde las
proporciones (un lugar el doble de lleno no saldría el doble).

\textbf{Paso 2, promediar con pesos iguales lo que el lugar tiene.}
\[
IPT_{d,t}=\frac{1}{|A_{d,t}|}\sum_{k\in A_{d,t}} z_{k,d,t},
\]
con $A_{d,t}$ = componentes con dato. Un componente entra al índice solo si lo tienen al menos dos lugares (si no, ese
lugar se compararía contra sí mismo; así salió Belice, que solo tiene Chetumal).

\textbf{Paso 3, clasificar con percentiles comunes.} $u_{50}$ y $u_{90}$ son los percentiles 50 y 90 del índice de todos
los lugares y meses:
\[
\text{estado}=\begin{cases}\text{tranquilo} & IPT<u_{50}\\ \text{concurrido} & u_{50}\le IPT\le u_{90}\\ \text{saturado} & IPT>u_{90}\end{cases}
\]
\emph{Fundamento}: un percentil es una posición relativa. ``Saturado'' significa ``en el 10\,\% más presionado de todo lo
observado en el estado'', la misma vara para el norte y el sur.

\begin{amano}[ · Cancún, julio de 2026 (888,797 habitantes)]
\begin{center}
\small
\begin{tabular}{lrrrr}
\encabezado \h{Componente} & \h{$x$} & \h{mín} & \h{máx} & \h{$z$} \\
Tren Maya (22,316 personas) & 25.1081 por mil & 0 & 505.2287 & 0.0497 \\
INAH, El Rey (854) & 0.9608 por mil & 0 & 6,451.8787 & 0.0001 \\
Llegadas por cuarto & 0.4709 & 0 & 432.0230 & 0.0011 \\
Ocupación (DataTur) & 68.78\,\% & 17.97 & 88.40 & 0.7214 \\
\end{tabular}
\normalsize
\end{center}
Por ejemplo: $z_{\text{tren}}=\dfrac{25.1081-0}{505.2287-0}=0.0497$ y $z_{\text{ocupación}}=\dfrac{68.78-17.97}{88.40-17.97}=\dfrac{50.81}{70.43}=0.7214$.
\[
IPT=\frac{0.0497+0.0001+0.0011+0.7214}{4}=\frac{0.7723}{4}=0.1931 .
\]
\end{amano}

\subsection{La prueba de validez que falló, y por qué se cuenta}
La primera versión (solo con SITUR-Q) calificaba a Cancún ``tranquilo'' en julio de 2026, con 0.025. La razón: SITUR-Q
dejó de publicar la ocupación en 2025, y sin ella Cancún solo tenía componentes casi en cero. Para ese mes, DataTur
marcaba 68.8\,\%. \textbf{Un índice que dice que Cancún está tranquilo en julio no es válido.} Se compararon cuatro
variantes con el IPT promedio de 2025–2026:

\begin{center}
\small
\begin{tabularx}{\linewidth}{>{\raggedright\arraybackslash}p{0.27\linewidth}rrrrY}
\encabezado \h{Variante} & \h{Cancún} & \h{Isla M.} & \h{Chetumal} & \h{Ruta} & \h{Problema} \\
A. Mín–máx, solo SITUR-Q & 0.02 & — & 0.39 & 0.14 & El norte ``desaparece'' \\
B. + ocupación DataTur & 0.28 & 0.45 & 0.39 & 0.14 & Chetumal arriba de Cancún \\
C. Percentil + DataTur & 0.44 & 0.25 & 0.48 & 0.71 & La Ruta arriba de Cancún \\
\textbf{D. B sin componentes de un solo lugar} & \textbf{0.28} & \textbf{0.45} & \textbf{0.05} & \textbf{0.14} & Quiebre de 2025 \\
\end{tabularx}
\normalsize
\end{center}
Brandon eligió \textbf{D}: la ocupación de los 4 lugares que mide DataTur viene de DataTur en toda su historia, y sale
Belice (solo Chetumal lo tiene).

\subsection{Dos índices, y por qué Cancún sale distinto en cada uno}
\begin{itemize}
\item \textbf{Índice base}: usa todas las medidas que existan cada mes. Cortes $u_{50}=0.258$, $u_{90}=0.766$. Con él,
  Cancún (0.193) es \emph{tranquilo}.
\item \textbf{El quiebre de 2025}: SITUR-Q dejó de publicar la ocupación del sur y Chetumal ``cayó'' de 0.54 a 0.05 sin que
  llegara menos gente. Fue un cambio en los datos, no en la realidad. Un modelo entrenado con eso aprendería cambios
  falsos.
\item \textbf{Índice comparable} (el que publica el Radar): cada lugar usa solo las medidas que tiene hoy, y solo en los
  meses en que las tiene todas. Así un cambio de estado es siempre real. Cortes $u_{50}=0.186$, $u_{90}=0.756$. Con él,
  Cancún (el mismo 0.193) es \emph{concurrido}.
\end{itemize}

\textbf{Julio de 2026, índice comparable}: concurridos Mahahual (0.731), Isla Mujeres (0.510), Bacalar (0.314), Holbox
(0.228), Cancún (0.193) y Cozumel (0.191); tranquilos Playa del Carmen (0.160), Tulum (0.145) y los cinco lugares de la
campaña (Ruta 0.129, Chetumal 0.025, Oxtankah 0.020, Maya Ka'an 0.009). La Laguna Milagros no tiene ninguna serie: ``sin
dato oficial'', nunca ``tranquila''.

\subsection{Sensibilidad: ¿y si los pesos no fueran iguales?}
Cada peso se movió a 0.5 y a 1.5 (con los demás en 1) y se contó cuántos lugar-mes cambian de estado: entre el 0.6\,\% y
el 6.0\,\%. Lo que más mueve el resultado es la ocupación. Conclusión: la elección de pesos iguales no decide el
resultado; el estado de casi todos los lugares es el mismo con cualquier peso razonable.

\begin{figure}[h]\centering
\includegraphics[width=0.85\linewidth]{f08_radar.png}
\caption{El Radar: índice comparable de cada lugar (notebook 02).}
\end{figure}

\begin{codigoen}\codigo{radar/panel.py} (panel), \codigo{radar/indice.py} (\codigo{minmax}, \codigo{ipt}, \codigo{estados},
\codigo{sensibilidad}), \codigo{radar/prediccion.py} (\codigo{indice_comparable}).\end{codigoen}

\begin{teprofe}
\emph{``¿Por qué pesos iguales?''} Porque no hay un dato que diga que el tren pesa más que los cruceros, y unos pesos
inventados serían difíciles de defender. Además, la sensibilidad muestra que con pesos de 0.5 a 1.5 el estado cambia en
6\,\% de los casos como máximo.
\end{teprofe}

\section{Agrupamiento de los centros turísticos del país (Ward)}
\textbf{Pregunta.} ¿A qué se parece el norte de Quintana Roo dentro del país? Si está en el grupo de los destinos más
llenos de México, la saturación no es una impresión.

\begin{intuicion}
Cada centro turístico tiene una ``huella'' de 12 números: qué tan lleno está cada mes del año. El agrupamiento junta los
centros cuyas huellas se parecen. Empieza con cada centro solo y va uniendo, de dos en dos, los grupos que menos
``desordenan'' el conjunto, hasta formar un árbol.
\end{intuicion}

\textbf{Fórmulas.} El perfil del centro $c$ es su ocupación de cada mes del año, sumando cuartos de todos los años, y la
distancia entre dos centros es la euclidiana:
\[
p_{c,m}=100\cdot\frac{\sum_{\text{años}}\text{ocupados}_{c,m}}{\sum_{\text{años}}\text{disponibles}_{c,m}},\qquad
d(c,c')=\sqrt{\textstyle\sum_m (p_{c,m}-p_{c',m})^2} .
\]
En cada paso, Ward une los grupos $A,B$ con el menor aumento de la varianza dentro de los grupos:
\[
\Delta(A,B)=\frac{|A|\,|B|}{|A|+|B|}\,\lVert \bar p_A-\bar p_B\rVert^2 .
\]
\emph{Fundamento}: la suma de cuadrados dentro de los grupos es lo que se quiere mínima (grupos compactos); unir $A$ y $B$
la aumenta exactamente en $\Delta(A,B)$, así que Ward elige en cada paso la unión más ``barata''.

\textbf{¿Cuántos grupos?} Con la silueta de cada centro, $s(i)=\dfrac{b(i)-a(i)}{\max\{a(i),b(i)\}}$, donde $a(i)$ es su
distancia media a los de su grupo y $b(i)$ al grupo más cercano. Va de $-1$ (mal asignado) a 1 (perfecto). Se elige el
número de grupos con la mayor silueta promedio:

\begin{center}
\begin{tabular}{lrrrrrrr}
\encabezado \h{$k$} & \h{2} & \h{3} & \h{4} & \h{5} & \h{6} & \h{7} & \h{8} \\
Silueta & \textbf{0.450} & 0.380 & 0.365 & 0.300 & 0.283 & 0.317 & 0.306 \\
\end{tabular}
\end{center}

\begin{amano}[ · silueta de Cancún]
$a=43.84$ (distancia media a los otros 22 de su grupo), $b=116.21$ (a los 32 del otro):
$s=\dfrac{116.21-43.84}{116.21}=\dfrac{72.37}{116.21}=0.623$. Cancún está bien asignado.
\end{amano}

\textbf{Resultado ($k=2$)}: grupo 1, playas y destinos muy ocupados (23 centros, 64.5\,\%: Cancún, Riviera Maya, Playacar,
Akumal, Playa del Carmen, Los Cabos); grupo 2, el resto del país (32 centros, 41.1\,\%, con Cozumel e Isla Mujeres). Solo
entran centros con sus 55 meses completos; los incompletos no se rellenan, se excluyen.

\begin{codigoen}\codigo{radar/clustering.py} (\codigo{perfiles}, \codigo{agrupar}): \codigo{scipy.cluster.hierarchy.linkage(method="ward")}
y \codigo{sklearn.metrics.silhouette_score}.\end{codigoen}

\section{La forma del año (estacionalidad)}
\label{sec:forma}
\textbf{Pregunta.} ¿Qué meses se llenan y cuáles se vacían en cada lugar? Es la base de la temporada alta, del pronóstico
y del calendario de la campaña.

\begin{intuicion}
Para saber si enero es ``alto'', se compara enero con el promedio de su propio año: si un año tuvo 5,000 visitantes al
mes en promedio y enero tuvo 8,000, enero fue 1.6 veces un mes normal. Se repite con varios años y se promedia.
\end{intuicion}

\subsection{Las fórmulas}
Descomposición clásica \textbf{multiplicativa}, solo con años completos:
\begin{align*}
r_{a,m}&=\frac{y_{a,m}}{\bar y_a},\qquad \bar y_a=\tfrac{1}{12}\sum_{m=1}^{12}y_{a,m} &&\text{(razón del mes contra su año)}\\
\tilde S_m&=\frac{1}{|A|}\sum_{a\in A}r_{a,m},\qquad S_m=\frac{\tilde S_m}{\frac{1}{12}\sum_k\tilde S_k} &&\text{(índice del mes, reescalado a promedio 1)}
\end{align*}

\textbf{Fundamento}: el modelo supone $y_{a,m}=\text{nivel}_a\times S_m\times\text{ruido}$. Al dividir entre el promedio
del año se elimina el nivel, y al promediar varios años se cancela el ruido. Queda la forma.

\textbf{¿Multiplicativa o aditiva?} Aditiva diría ``enero tiene 2,000 visitantes más''; multiplicativa, ``enero tiene 61\,\%
más''. Kohunlich reabrió en 2025 con la mitad de su nivel de 2019: ``+2,000'' ya no vale, pero ``+61\,\%'' sí. La temporada
escala con el nivel, así que es multiplicativa.

\textbf{Fuerza de la temporada} (Wang, Smith y Hyndman, 2006), en logaritmos, con $s=\ln S_m$ y $e=\ln r_{a,m}-\ln S_m$:
\[
F_S=\max\!\left(0,\;1-\frac{\Var(e)}{\Var(s+e)}\right).
\]
\emph{Fundamento}: $s+e$ es toda la variación de un mes a otro; $e$ es la parte que la temporada no explica. Si la
temporada lo explica todo, $\Var(e)=0$ y $F_S=1$; si no explica nada, $F_S=0$.

\begin{amano}[ · Ruta arqueológica del sur]
\textbf{Una razón.} En 2019, Kohunlich + Dzibanché sumaron 64,139 visitantes: $\bar y_{2019}=64{,}139/12=5{,}344.92$. En
enero hubo 7,935: $r_{2019,1}=7{,}935/5{,}344.92=1.4846$.

\textbf{El índice de enero.} Las seis razones de enero (2016–2019, 2022, 2023): 1.5459, 1.7075, 1.8266, 1.4846, 1.2540 y
1.8173. Suman 9.6359: $\tilde S_1=9.6359/6=1.606$. El promedio de los 12 índices es 1.0000, así que $S_1=1.606$: en enero
llega 61\,\% más gente que en un mes promedio. En septiembre, $S_9=0.496$: la mitad.

\textbf{La fuerza.} $\Var(e)=0.03152$ y $\Var(s+e)=0.15230$:
\[F_S=1-\frac{0.03152}{0.15230}=1-0.2070=0.793 .\]
El 79\,\% de la variación de un mes a otro en la Ruta es temporada.
\end{amano}

\begin{center}
\begin{tabular}{lrll}
\encabezado \h{Serie} & \h{$F_S$} & \h{Mes más alto} & \h{Mes más bajo} \\
Ruta arqueológica del sur & 0.793 & enero (1.606) & septiembre (0.496) \\
Bahía Calderitas–Oxtankah & 0.684 & diciembre (1.402) & septiembre (0.647) \\
Chetumal (cruces de Belice) & 0.655 & diciembre (1.172) & febrero (0.865) \\
Cancún (referencia) & 0.714 & marzo (1.084) & septiembre (0.867) \\
\end{tabular}
\end{center}

\textbf{¿Por qué no STL?} STL (descomposición con regresión local, el método del plan) necesita una serie continua, y estas
tienen huecos: pandemia, cierres por obras. STL solo podría usar el tramo continuo más largo y tiraría los años
posteriores. Se usó como \textbf{segunda opinión} en ese tramo: su índice coincide con el de este método entre 0.85 y
0.97.

\begin{figure}[h]\centering
\includegraphics[width=0.85\linewidth]{f10_forma_del_anio.png}
\caption{La forma del año de cada serie (notebook 03).}
\end{figure}

\begin{codigoen}\codigo{pronostico/forma.py}: \codigo{razones}, \codigo{indice_estacional}, \codigo{fuerza_estacional},
\codigo{segunda_opinion_stl}.\end{codigoen}

\section{Clima raro con Isolation Forest (detección de anomalías)}
\label{sec:iforest}
\textbf{Pregunta.} ¿Esta semana tuvo un clima tan raro que conviene pausar el anuncio? Se necesita detectar ``lo raro''
sin que alguien defina a mano qué es raro.

\begin{intuicion}
Imagina que separas puntos en un plano con cortes al azar. Un punto que está lejos de todos se queda solo con muy pocos
cortes; un punto en medio de la multitud necesita muchísimos. Isolation Forest hace eso con cientos de árboles de cortes
al azar y mide cuántos cortes tarda, en promedio, en aislar cada semana. Las semanas raras se aíslan rápido.
\end{intuicion}

\subsection{Las fórmulas y su fundamento}
\[
s(x)=2^{-\frac{E[h(x)]}{c(n)}},\qquad c(n)=2H(n-1)-\frac{2(n-1)}{n},\qquad H(k)\approx \ln k+0.5772 .
\]
\begin{itemize}
\item $h(x)$: número de cortes hasta aislar la semana $x$ en un árbol; $E[h(x)]$, su promedio en los 300 árboles.
\item $c(n)$: el número \emph{promedio} de cortes que necesita una semana cualquiera en un árbol de $n$ datos. Viene de la
  teoría de árboles binarios de búsqueda: es la longitud media de una búsqueda sin éxito. Sirve para normalizar.
\item Si $E[h(x)]\approx c(n)$, entonces $s\approx 2^{-1}=0.5$ (normal). Si $E[h(x)]\ll c(n)$, $s\to1$ (muy rara).
\item Cada árbol se entrena con una submuestra de hasta 256 semanas: con pocas semanas por árbol, las anomalías se
  separan mejor (el artículo original lo demuestra).
\end{itemize}

\textbf{Variables}: lluvia de la semana, lluvia del peor día, viento máximo, temperatura máxima y la época del año como seno
y coseno de la semana. La época entra para que un aguacero de septiembre no salga raro solo por ser septiembre.

\begin{amano}[ · Chetumal, semana del 14 de octubre de 2024 (tormenta Nadine)]
Llovieron 148.7 mm (65.4 mm el peor día). Había 260 semanas de entrenamiento, así que cada árbol usa $n=256$:
\[
c(256)=2(\ln 255+0.5772)-\frac{2\cdot255}{256}=2(5.5413+0.5772)-1.9922=12.237-1.992=10.245 .
\]
El bosque da $s=0.712$. El corte (percentil 95 de las semanas de entrenamiento) es $0.543$: como $0.712>0.543$, la semana
es \textbf{rara}. ¿Cuántos cortes tardó en aislarse?
\[
E[h]=-c(n)\log_2 s=-10.245\times\log_2 0.712=-10.245\times(-0.489)=5.0 ,
\]
contra unos 10 de una semana normal.
\end{amano}

\begin{decision}
El corte automático del método (\codigo{contamination="auto"}, que marca raro si $s>0.5$) señalaba \textbf{47\,\%} de las
semanas de Chetumal: con pocas semanas el bosque no distinguía, y desde 2022 llueve más y hace más calor que en 2019–2021.
Brandon eligió ``más raro que el 95\,\% de las semanas que el bosque ya conoce'', con un bosque por año entrenado con
todos los años anteriores (ventana creciente). Resultado: 21 semanas raras en Chetumal y 17 en Kohunlich.
Descartados: el percentil 99 (casi nunca pausaba) y quitar la señal.
\end{decision}

\textbf{Validación independiente}: el bosque marcó como raras las tres semanas en que pasaron tormentas (Lisa en 2022,
Nadine y Sara en 2024) \textbf{sin saber} que hubo tormenta. Es evidencia de que ``raro'' significa algo real.

\begin{codigoen}\codigo{envivo/senales.py}: \codigo{clima_semanal}, \codigo{anomalias_clima}
(\codigo{sklearn.ensemble.IsolationForest(n_estimators=300, random_state=0)}).\end{codigoen}

\section{Minería de texto de las reseñas}
\label{sec:texto}
\textbf{Pregunta.} ¿Qué molesta y qué enamora al turista? Es lo que define la promesa de la campaña y sus palabras.

\textbf{Datos}: 85,987 reseñas de Quintana Roo (Rest-Mex 2025, sin duplicados): Tulum, Isla Mujeres y Bacalar.
\textbf{Ninguna es de los cinco lugares}, así que sirven para saber qué molesta en los destinos llenos, no para opinar
del sur. 3,597 son malas (1–2 estrellas) y 60,082 de 5 estrellas.

\subsection{Temas con un léxico a la vista y riesgo relativo}
Cada tema es una lista de raíces de palabras que cuentan solo al inicio de una palabra. Por ejemplo, ``ruido'' incluye
\emph{ruido, ruidos, escándalo, música alta, fiesta}; ``calma'' incluye \emph{tranquil-, paz, relaj-, silencio}.

El \textbf{riesgo relativo} es la medida que usa la epidemiología para comparar el riesgo en un grupo con el riesgo
general:
\[
RR_a=\frac{P(\text{mala}\mid\text{la reseña toca el tema }a)}{P(\text{mala})}.
\]
$RR=1$: el tema no cambia nada. $RR=2$: una reseña que lo toca es mala el doble de seguido.

\begin{amano}[ · ruido]
$P(\text{mala})=3{,}597/85{,}987=0.04183$. Reseñas que hablan de ruido: 2,038, de las cuales 235 son malas:
$P(\text{mala}\mid\text{ruido})=235/2{,}038=0.11531$. Entonces
\[RR_{\text{ruido}}=\frac{0.11531}{0.04183}=2.76 .\]
\end{amano}

\begin{center}
\begin{tabular}{lrl}
\encabezado \h{Tema} & \h{$RR$} & \h{Lectura} \\
Ruido & 2.76 & Aleja \\ Suciedad & 1.89 & Aleja \\ Precio & 1.66 & Aleja \\ Multitudes & 1.52 & Aleja \\
Servicio & 1.18 & Neutro \\ Comida & 0.81 & Protege un poco \\ Naturaleza & 0.73 & Protege \\
Cultura & 0.48 & Protege \\ Calma & 0.41 & Protege \\
\end{tabular}
\end{center}

\textbf{Por qué léxico y no un modelo de temas (LDA)}: con LDA, cada tema es una lista de palabras que hay que interpretar
(``adivinar'' qué significa). Con un léxico, cada tema se define antes y se defiende palabra por palabra. Para una campaña
que tiene que explicar por qué promete calma, lo segundo es más defendible.

\subsection{Reglas de asociación (algoritmo Apriori)}
Cada reseña es una ``canasta'' con sus temas y su calificación. Apriori encuentra combinaciones frecuentes:
\[
\text{soporte}(A)=\frac{\#\{\text{reseñas con }A\}}{N},\qquad
\text{confianza}(A\Rightarrow B)=\frac{\text{sop}(A\cup B)}{\text{sop}(A)},\qquad
\text{lift}=\frac{\text{confianza}}{\text{sop}(B)} .
\]
\emph{El truco de Apriori}: si una combinación es rara, todas las combinaciones más grandes que la contienen también lo son
(propiedad antimonótona). Así no tiene que revisar todas las combinaciones posibles. Se usó soporte mínimo de 0.1\,\% y
combinaciones de hasta 3 elementos.

\begin{amano}[ · ruido + servicio $\Rightarrow$ reseña mala]
1,204 reseñas tocan ruido y servicio a la vez; 142 de ellas son malas.
\[
\text{soporte}=\frac{142}{85{,}987}=0.00165,\quad \text{confianza}=\frac{142}{1{,}204}=0.1179,\quad
\text{lift}=\frac{0.1179}{0.04183}=2.82 .
\]
Una reseña que habla de ruido y de mal servicio es mala 2.8 veces más seguido de lo normal. Otras: suciedad + multitudes
(lift 2.59), suciedad + precio (2.54). Hacia 5 estrellas: comida + servicio (1.12), calma + naturaleza (1.10).
\end{amano}

\subsection{El lenguaje del turista (log-odds con prior de Dirichlet)}
\textbf{Pregunta}: ¿qué palabras usan más las reseñas de 5 estrellas que las malas? Contar no basta: las palabras raras dan
cocientes enormes por azar. El método de Monroe, Colaresi y Quinn (2008) agrega un ``colchón'' $\alpha_w$ proporcional a lo
común que es la palabra:
\[
\delta_w=\ln\frac{y^A_w+\alpha_w}{n_A+\alpha_0-y^A_w-\alpha_w}-\ln\frac{y^B_w+\alpha_w}{n_B+\alpha_0-y^B_w-\alpha_w},
\qquad z_w=\frac{\delta_w}{\sqrt{\frac{1}{y^A_w+\alpha_w}+\frac{1}{y^B_w+\alpha_w}}} ,
\]
con $A$ = reseñas de 5 estrellas, $B$ = reseñas malas, $y$ = veces que aparece la palabra, $n$ = total de palabras,
$\alpha_w=0.01\times$(veces en todo el corpus) y $\alpha_0=\sum_w\alpha_w$. $z$ es como un puntaje de cuántas
desviaciones estándar se aleja la palabra de lo esperado.

\begin{amano}[ · la palabra ``excelente'']
$y^A=18{,}952$, $y^B=144$, $\alpha=190.96$, $n_A=1{,}879{,}868$, $n_B=162{,}035$, $\alpha_0=20{,}419.0$.
\begin{align*}
\delta&=\ln\frac{18{,}952+190.96}{1{,}879{,}868+20{,}419.0-18{,}952-190.96}-\ln\frac{144+190.96}{162{,}035+20{,}419.0-144-190.96}\\
&=\ln\frac{19{,}142.96}{1{,}881{,}144.04}-\ln\frac{334.96}{182{,}119.04}=-4.5878-(-6.2985)=1.7107,\\
z&=\frac{1.7107}{\sqrt{1/19{,}142.96+1/334.96}}=\frac{1.7107}{0.05512}=31.04 .
\end{align*}
``Excelente'' es una de las palabras que más distinguen a las reseñas de 5 estrellas. En el otro extremo, ``sucia'' tiene
$z=-18.76$.
\end{amano}

Palabras de 5 estrellas: \emph{excelente, increíble, hermoso, deliciosa, amable, ruinas, ambiente}. Malas: \emph{peor,
horrible, dinero, caro, sucia, grosero}. Con ellas se escribieron los mensajes (capítulo~\ref{cap:mkt}).

\begin{codigoen}\codigo{campana/texto.py}: \codigo{marcar_aspectos}, \codigo{aspectos}, \codigo{reglas_asociacion}
(\codigo{mlxtend.apriori}), \codigo{palabras_que_distinguen}. Notebook \codigo{07_campana}.\end{codigoen}

\section{NLP de lugares: qué hacer, dónde comer y dónde dormir}
\textbf{Pregunta}: de los negocios del DENUE, ¿cuáles recomendar de día, tarde y noche en cada lugar?

\textbf{Método}: un clasificador por léxico que lee el \textbf{nombre} del negocio y, si el nombre no dice nada, su giro
oficial (SCIAN). No se entrenó un modelo estadístico porque no hay etiquetas con qué entrenarlo; un léxico es lo honesto y
se puede explicar regla por regla. Las reglas salieron de errores reales:
\begin{itemize}
\item La raíz cuenta solo al inicio de una palabra: ``MICHELADAS'' contiene ``HELAD'' y salía como heladería;
  ``BARBACOA'' contiene ``BAR''.
\item El nombre gana al giro: ``HOTEL COSTA AZUL'' tenía giro de casa de huéspedes.
\item Se excluye lo que no le sirve a un visitante: cooperativas escolares, gimnasios, moteles, banquetes.
\end{itemize}
\textbf{Precisión medida}: dos muestras al azar revisadas a mano, 38 de 40 y, tras corregir, \textbf{39 de 40} en una
muestra nueva. \textbf{No se hizo scraping} de Google Maps (lo prohíben sus términos): cada tarjeta es un enlace público.

\begin{codigoen}\codigo{campana/lugares.py}: \codigo{clasificar}, \codigo{recomendaciones}.\end{codigoen}
